\hypertarget{appendix-1f9963e1a5d1}{}
\section*{Bit packing}
\section{Bit packing}\label{appendix-1f9963e1a5d1-bit-packing}

\textbf{Curriculum depth:} Introduced · \textbf{First seen:} C001 · \textbf{Latest development:} C001

\subsection{Definition and mechanism}\label{appendix-1f9963e1a5d1-definition-and-mechanism}

Multiple constrained values share a carrier word; extraction and mutation must preserve the declared schema.

Packing places several constrained values into a carrier. Input \texttt{{[}1,0,1,1,0,1,0,1{]}} packed LSB-first produces \texttt{0xAD}. Unpacking needs the logical count because the final byte may contain unused bits. Define whether padding must be zero, ignored, or rejected.

\subsection{Prerequisites}\label{appendix-1f9963e1a5d1-prerequisites}

\begin{itemize}
\tightlist
\item
  \hyperref[appendix-ac697f5c66e8]{Bit field} --- A field is a declared range with offset s and width w, not automatically a C++ language bit-field member.
\item
  \hyperref[appendix-a427ce8ce8a6]{Bitwise OR} --- Per lane: only 0\textbar0 is 0; the other three cases are 1.
\item
  \hyperref[appendix-f414bd4ed05c]{Interpretation contract} --- For \texttt{0xAD}, unsigned interpretation gives 173; eight-bit two's complement gives −83.
\end{itemize}

\subsection{Used by}\label{appendix-1f9963e1a5d1-used-by}

\begin{itemize}
\item
  \hyperref[appendix-638204c40e8b]{Groupwise quantization and metadata} --- uses this mechanism to define its representation, error or validation contract.
\item
  \hyperref[appendix-e12d1b03fefb]{Atomic operations and write contention} --- uses bit packing as part of its representation or reasoning.
\item
  \href{../04-Visual-Computing-II/main.pdf\#nameddest=appendix-fdc7879e8bf2}{Quantization and compressed weights (04-Visual-Computing-II, appendix-fdc7879e8bf2)} --- uses bit packing as part of its representation or reasoning.
\end{itemize}

\subsection{Recall and next depth}\label{appendix-1f9963e1a5d1-recall-and-next-depth}

Explain bit packing without the label, work the example above, then state the assumption whose removal would change the result.

\textbf{Pending:} bulk kernels, padding policy, random access, atomic mutation, and format evolution.

\subsection{Source and navigation}\label{appendix-1f9963e1a5d1-source-and-navigation}

\hyperref[01-foundations-i-ch001-from-a-signal-to-an-interpretation]{C001 teaching chapter} · \hyperref[appendix-fd090d982a15]{Part definitions and mechanisms} · \href{../02-Foundations-II/main.pdf\#nameddest=appendix-4ad47bdd269a}{Dependency Map (02-Foundations-II, appendix-4ad47bdd269a)}
